\section{Finite ternary phase router and a positive binary $H_2$ word}
\label{app:finite-phase}

The scalar bank-exchange and frame-label architecture follows the
finite-network method of~\cite{swap}; its characteristic-two and
ternary-phase adaptation is made explicit below. These finite-network
antecedents are credited independently of the present tape-time claim.

This appendix records the finite algebraic and rank-saving part of the
binary $H_2$ supplier. It distinguishes three coefficient/address domains:
the address-frame geometry is over $\mathbb F_3$, the phase kernel and
dirty-selector payload arithmetic are over $\mathbb F_4$, and the fixed
positive word is over its codec field $\mathbb F_{16}$. The tape costs for
the equal-chunk interchange and weighted native mask are separate machine
suppliers; no field element in these algebraic arguments is treated as a
unit-cost operation on a growing coefficient word.

\subsection{The ternary phase kernel}

Let $\eta\in\mathbb F_4$ satisfy $\eta^2+\eta+1=0$. On $\mathbb F_4^3$ put
\[
 C=\eta I_3+\eta^2J_3
 =\begin{pmatrix}1&\eta^2&\eta^2\\
 \eta^2&1&\eta^2\\\eta^2&\eta^2&1\end{pmatrix}.
\]
Then $C^3=I$. Let $\chi(a)=\eta^a$ for $a\in\mathbb F_3$, and let
\[
 (\mathcal F_3)_{x,z}=\chi(xz),\qquad x,z\in\mathbb F_3.
\]
Since $1+\eta^a+\eta^{2a}=0$ for nonzero $a\in\mathbb F_3$,
$\mathcal F_3^2$ is the coordinate-negation permutation. Thus
$\mathcal F_3^4=I$ and $\mathcal F_3^{-1}=\mathcal F_3^3$ (equivalently,
$(\mathcal F_3^{-1})_{x,z}=\eta^{-xz}$).
Direct multiplication gives
\[
 C=\mathcal F_3\,\operatorname{diag}(1,\eta,\eta)\,\mathcal F_3^{-1}.
\]
Thus the eigenphase of $C$ at frequency $z$ is $\eta^{z^2}$. The inverse
is $C^2$ and has eigenphase $\eta^{-z^2}$.

For an address space with a nondegenerate symmetric form over $\mathbb F_3$,
write $T_v f(x)=f(x+v)$. If $v$ has norm $\nu=\langle v,v\rangle\in
\{1,2\}$, define
\[
 K_v=I+\eta^{-\nu}(T_v+T_{-v}).                         \tag{A.1}
\]
On a Fourier character of frequency $z$, the two translations have
eigenvalues $\eta^{\langle v,z\rangle}$ and
$\eta^{-\langle v,z\rangle}$. For $a=\langle v,z\rangle\in\mathbb F_3$,
\[
 1+\eta^{-\nu}(\eta^a+\eta^{-a})=\eta^{a^2/\nu}.     \tag{A.2}
\]
Indeed both sides are $1$ for $a=0$; for $a=\pm1$, checking
$(\nu,a^2)=(1,1)$ and $(2,1)$ gives $\eta$ and $\eta^2$ respectively.
Consequently $K_v$ is precisely the one-direction quadratic phase for
the line $\langle v\rangle$. Orthogonal directions commute, and an
orthogonal direct sum of directions has the product of their phases.
Replacing $v$ by $-v$ does not change the line; norm $2$ replaces $C$ by
$C^{-1}$.

For a nondegenerate subspace $U$ with an orthogonal basis
$v_1,\ldots,v_r$, define $C_U=\prod_jK_{v_j}$. In Fourier coordinates
this is the multiplier
\[
 (\mathcal F C_U\mathcal F^{-1})_{z,z}
   =\eta^{\sum_j\langle v_j,z\rangle^2/\langle v_j,v_j\rangle}.
                                                               \tag{A.3a}
\]
Thus orthogonal direct sums multiply their $C_U$ operators. If
$U\subseteq V$ and $V=U\perp W$, then $C_V=C_UC_W$; a rank-$r$
edge between such frames needs exactly $r$ one-direction children.
The phase frame on a role is a proof invariant: the same $\mathbb F_4$-linear
phase operator is applied to aligned wire arrays, so it commutes with
pointwise XOR/addition. It is a linear operator on payload arrays, not an
address relabeling; this is distinct from the address permutations used
in the SWAP supplier.
At completion the role permutation pairs each input frame with its
specified output frame. Choosing the X-source inverse line frame makes
the endpoint product the full phase, while its first edge has difference
$P_U-(-P_U)=2P_U$ and therefore one rank-one child. This is the source
correction counted below.

\subsection{The $48$-triple frame over $\mathbb F_3$}

Fix $H=48$, let $\mathcal T=\binom{[H]}3$, and join distinct $S,T$ when
$|S\cap T|=1$. Put
\[
 v=\binom{48}{3}=17296,\quad
 z=3\binom{45}{2}=2970,\quad N=v^3,\quad I=3v^2.
\]
The factor space is $F=\mathbb F_3^{49}$ with form
\[
 B=\operatorname{diag}(-I_{48},1).
\]
For $S\in\mathcal T$ use $t_S=(1_S,1)$. Then
\[
 \langle t_S,t_S\rangle=-3+1=-2=1\pmod 3,
 \qquad \langle t_S,t_T\rangle=1-|S\cap T|.
\]
Thus the line $L_S=\langle t_S\rangle$ is unit norm and adjacent triples
give orthogonal lines. The form is nondegenerate. To identify it with
standard dot-product coordinates, split the first $48$ coordinates into
$24$ disjoint pairs, and on each pair use
\[
 (a,b)\longmapsto(2a+2b,\,2a+b).
\]
The two output vectors for the basis inputs are $(2,2)$ and $(2,1)$;
their norms are $2$ and $2$, and their dot product is $2\cdot2+2\cdot1
=0$ in $\mathbb F_3$. The map is invertible (determinant $1$ in $\mathbb F_3$), so this is
an isometry from $\operatorname{diag}(-1,-1)$ to the standard dot
product. The last coordinate already has norm one.

For a stage prefix tensor space $A=F^{\otimes(j-1)}$, let $P$ be its
prefix line, $B_0=P^\perp$, and let $L_t$ denote the current triple line.
Since $P$ and $L_t$ are nondegenerate,
\[
 A=B_0\perp P,\qquad F=L_t\perp L_t^\perp,
\]
and
\[
 (P\otimes L_t)^\perp
   =(B_0\otimes F)\perp(P\otimes L_t^\perp).        \tag{A.3}
\]
The neighbor condition supplies $L_s\perp L_t$ whenever a side role
contains the ordered neighbor $(s,t)$. These decompositions prove the
following eight gate labels are nested/comparable exactly where their
wire edges require:
\[
\begin{array}{c|l}
 t_0& B_0\otimes L_t\\
 t_1& B_0\otimes F\\
 t_2& (B_0\otimes F)\perp(P\otimes L_t)\\
 t_3& A\otimes F\\
 t_4& B_0\otimes F\\
 t_5& (B_0\otimes F)\perp(P\otimes L_t^\perp)\\
 t_6& A\otimes F\\
 t_7& A\otimes F.
\end{array}                                                   \tag{A.4}
\]
In stage two the physical wire order is the inverse of the forward row
schedule with logical source and target exchanged; this keeps the same
physical X/Y label convention in (A.4). In particular, its incoming
physical X label is $(A\otimes L_t)\otimes Q$ and its outgoing label
$(A\otimes F)\otimes Q$; its incoming physical Y label is
$(B_0\otimes L_t)\otimes Q$ and outgoing label
$(P\otimes L_t)^\perp\otimes Q$. The latter is (A.3) for the next
stage, with the suffix tensor line $Q$ restored. This checks the
interstage edges as well as the within-invocation inclusions.

Here is the wire-level interpretation of the table. A stage fixes the
other two triple indices of a data role in $\mathcal T^3$; all displayed
labels are tensored with its fixed suffix line $Q$. Name an ordered
neighbor pair by its physical Y and physical X triple indices in
that stage, renaming the ordered pairs in the inverse second stage.
The side wire
$a_{(S,T)}$ meets, in order, the Y-side gate for $S$, the X-side
gate for $T$, the Y-side gate for $S$, and the X-side gate for $T$.
Their labels are rows $t_0,t_2,t_5,t_7$. Its middle inclusion is
precisely $L_T\subset L_S^\perp$, supplied by the neighbor condition.
Each center wire meets $t_1,t_3,t_4,t_6$, whose labels alternate
$B_0\otimes F$, $A\otimes F$, $B_0\otimes F$, $A\otimes F$.
Data X wires meet $t_2,t_3,t_6,t_7$ and data Y wires meet
$t_0,t_1,t_4,t_5$. These descriptions check every gate incidence;
multi-input sums can be decomposed into the same common-label XOR
gates without adding frame edges of nonzero rank.
Side and center scratch wires have identity input phase and full-space
output phase, connecting their first/last displayed labels to zero/full
labels. Initial Y data have identity phase. Final X data have full-space
phase and final Y data have the complement of their triple tensor line.
The initial X inverse line is the sole exceptional source frame.
The physical scratch rows also require the full phase: they were
arbitrary rows of the original array. It is their logical scalar
scratch values that the routing motif restores before the frame is
interpreted at the output.

For a fixed eligible mask sector, let $P$ be this scalar role permutation
acting on aligned arrays, and let $D_{\rm in},D_{\rm out}$ be the
block-diagonal input/output phase frames. A frame change on an edge
from $F$ to $G$ applies $GF^{-1}$. Common-label XOR gates commute with
their shared phase, so the entire framed program telescopes to
$D_{\rm out}P D_{\rm in}^{-1}$. If input role $r$ is sent to $P(r)$,
the endpoint operator there is
$C_{{\rm out},P(r)}C_{{\rm in},r}^{-1}=C_{\rm full}$ for data
\emph{and} scratch roles. Consequently
\[
 D_{\rm out}P D_{\rm in}^{-1}
     =P(I_W\otimes C_{\rm full}).                       \tag{A.4a}
 \label{eq:full-role-phase}
\]
The inverse role merge removes $P$, leaving the full phase on every
arbitrary input role array. Eligibility is the same mask-only predicate
on every row role; in an ineligible sector every frame is identity,
so the scalar $P$ is likewise canceled by that merge. No scratch
payload is initialized to zero in this identity.

\subsection{Arbitrary-scratch scalar routing and the finite rank sum}

For one invocation over $\mathbb F_2$, let $x$ be its data source, $y$
its target, $a$ its side scratch vector and $c$ its center scratch vector.
Let $V,J,G,R$ be the incidence maps in the companion scalar-motif
certificate. The chronological eight rows are
\[
\begin{aligned}
 y&\gets y-Ja, & y&\gets y-Rc,\\
 a&\gets a+Vx, & c&\gets c+Gx,\\
 y&\gets y+Rc, & y&\gets y+Ja,\\
 c&\gets c-Gx, & a&\gets a-Vx.
\end{aligned}
\]
where in the first two rows $a,c$ are their entry values. In
characteristic two subtraction equals addition, but the signs make the
scratch cancellation visible. The final two rows restore the center and
side scratch. Equivalently,
for arbitrary initial scratch $(a_0,c_0)$ the total target increment is
\[
 -Ja_0-Rc_0+R(c_0+Gx)+J(a_0+Vx)=(RG+JV)x.             \tag{A.5}
\]
For triples, $G_i,T=1_{i\in T}$, $R_{S,i}=1_{i\in S}$,
$V_{(S,T),T'}=1_{T=T'}1_{|S\cap T|=1}$ and
$J_{S,(S,T)}=1_{|S\cap T|=1}$. If $S=T$, $RG=3=1$ in $\mathbb F_2$
and $JV=0$. If $S\ne T$ with intersection one, both terms are one
and cancel; for intersections zero or two both vanish. Hence
$RG+JV=I$. Formula (A.5) proves exact addition and restoration for every
scratch assignment.

Run the forward schedule X-to-Y, its inverse schedule Y-to-X, then the
forward schedule again. On arbitrary banks
\[
 (x,y)\mapsto(x,y+x)\mapsto(-y,y+x)\mapsto(-y,x),
                                                               \tag{A.6}
\]
so in characteristic two this exchanges the banks. Every invocation has
disjoint scratch roles; all are fixed by the completed role permutation.
There are $v$ roles per bank, $vz$ ordered side roles and $H$ center
roles per invocation, hence
\[
 W=2N+I(vz+48).                                           \tag{A.7}
\]

For a label of dimension $r$ over $\mathbb F_3$, a wire operation on
that label uses its quadratic phase and a source/target frame change of
rank equal to the relevant dimension difference. The only downward
label transitions in (A.4) are the $t_3\to t_4$ center transitions.
Each loses $F=49$ dimensions and there are $48$ center coordinates in
each of $I$ invocations. Their total downward loss is
\[
 L=I\cdot48\cdot49.                                    \tag{A.8}
\]
The signed dimension telescope over all routed wire endpoints is
$WM-2N$, where $M=49^3=117649$: the two banks contribute the two $N$
endpoint deficits. Turning signed variation into absolute variation
adds twice every downward loss, so before the source correction the sum
of edge ranks is $WM-2N+2L$.

This count can also be read from the complete physical wire paths.
Put $d=49$ and $a=d^{j-1}$ at stage $j$. With the suffix line
suppressed, the side and center dimensions are respectively
\[
 0,\ a-1,\ d(a-1)+1,\ da-1,\ da,\ M,
 \qquad
 0,\ d(a-1),\ da,\ d(a-1),\ da,\ M .
\]
The side inclusions follow from $L_T\subset L_S^\perp$, with
nondegenerate rank-$(d-2)$ residual. The center loses exactly
$d$ dimensions at $t_3\to t_4$ and otherwise increases.
Physical X enters the next stage as $A\otimes L_T\subset t_2(T)$;
physical Y enters as $B_0\otimes L_S=t_0(S)$.
Their complete ordinary paths have no decreases and have total
rank $M-1$ each, using a formal positive X input line for this count.
The inverse middle stage retains this physical order: its inverse
logical rows $7,\ldots,0$, after bank exchange, have the same physical
incidences as chronological rows $0,\ldots,7$. The label table itself
is not reversed. Consequently the total ordinary rank is
\[
 2N(M-1)+(W-2N)M+2I\cdot48d=WM-2N+2L.
\]

Choose the X-source phase frame to be the inverse line phase on $U_a$;
all other source frames are identity and gate/sink frames are the
positive phases of (A.4). The endpoint identity on each X-to-Y route is
$C_{U_a^\perp}C_{U_a}=C_{F^{\otimes3}}$. On each of the $N$ first X
edges, the source label is $-P_{U_a}$ while the first gate label is
$+P_{U_a}$, so their difference is $2P_{U_a}$. This has rank one and
its phase is $C_{U_a}^{-1}$, since scaling the quadratic form by
$2=-1$ in $\mathbb F_3$ inverts its phase. These are $N$ genuine extra
children, not a rank-zero equal-dimension edge. Therefore
\[
 s=WM-2N+2L+N=WM-N+2L.                                \tag{A.9}
\]

At $H=48$ direct cancellation gives
\[
 \frac WN=\frac{9633881}{1081},\qquad
 \frac LN=\frac{441}{1081},\qquad
 1-\frac{s}{WM}=\frac{199}{9633881\cdot117649}.
                                                               \tag{A.10}
\]
The integer comparison
\[
 199\,2^{36}>12\cdot9633881\cdot117649
\]
and $\log_2 M<17$, $\ln2<7/10$ imply $\ln M<12$ and
\[
 \frac{199}{9633881\cdot117649}>12\,2^{-36}
   >1-M^{-2^{-36}}.
\]
Thus $s/W<M^{1-2^{-36}}$. This is the explicit strict role saving used
below; no numerical matrix-rank test is involved.

\subsection{From the role word to all active sectors}

An $A$-digit has states $0,1,2$ active and state $3$ idle. Represent
$d_i=2X_i+Y_i$. For any active mask $S$, $A^{\otimes h}$ restricts to
$C^{\otimes |S|}$ on the active coordinates and identity on idle
coordinates. All $4^h$ addresses remain in the stream; no sector is
gathered out.

The mask-routing identities use dirty address fields. Let $T_Z$ apply
one fixed transposition of active symbols where dirty selector bit
$Z_i=1$. It preserves the active mask $S$, but arbitrary dirty $Z$
can make it change paired sources as well as targets. Source-value
predicates are recomputed only after the completed $T_{J\wedge Q}$
word below, which changes matched targets only. For a mask-only selector
$J(S)$ the chronological word
\[
 Z\mathrel{\oplus}=J(S),\quad T_Z,\quad
 Z\mathrel{\oplus}=J(S),\quad T_Z                         \tag{A.11}
\]
restores arbitrary initial $Z$ and equals $T_J$. With a second dirty
field $Q$, use $J(S)\wedge Q$ in (A.11). If $F_a(d)_i$ says the paired
source of target $i$ currently has trit $a$, then the completed
$T_{J\wedge Q}$ changes targets only and preserves all paired sources.
Consequently
\[
 Q\mathrel{\oplus}=F_a(d),\quad T_{J\wedge Q},\quad
 Q\mathrel{\oplus}=F_a(d),\quad T_{J\wedge Q}              \tag{A.12}
\]
restores arbitrary $Q,Z$ and implements the target transposition
conditioned on source value $a$. The source/target matching must be
disjoint; this is what makes the second $F_a$ equal to the first at the
completed inner-operation boundary. A third dirty field is reserved for
the weighted tile auxiliary. The supplied weighted-mask lemma implements
each frozen XOR in (A.11)--(A.12) with its separate auxiliary restored;
its exact sparse boundary repair and interchange cost are not reproved
here.

The two active transpositions have explicit bit circuits
\[
 (0\ 1):\ Y\mathrel{\oplus}=Z\wedge\neg X,
 \qquad
 (1\ 2):\ X\mathrel{\oplus}=Z\wedge Y;\quad
 Y\mathrel{\oplus}=Z\wedge X;\quad
 X\mathrel{\oplus}=Z\wedge Y.                           \tag{A.13}
\]
They are constant sequences of frozen-control XOR shears. Applying the
two transpositions in chronological order $(1\ 2)$ then $(0\ 1)$ adds
one to an active trit; reverse order subtracts one. Conditioning on
source trit and composing these shears realizes a matching SUM
\[
 \text{target}_i\gets\text{target}_i+\text{source}_i\pmod3.
                                                               \tag{A.14}
\]
Three SUMs plus a negation implement a coordinate swap. All dirty
selectors and the tile auxiliary return to their arbitrary incoming
values after each completed operation.

For a node acting on the first $e$ active ranks, restrict controls and
endpoints to the first
\[
 u=\min\{h,\lceil C_0(e+\log(L+2))\rceil\}
\]
physical positions, where $L\ge h$ is the ambient address-width bound
in the native implementation. If fewer than $e$ active positions occur there,
define the operation to be identity; otherwise all required ranks are
present. Restrict every raw $T_Z$ transposition and both selector-bank
toggles to this prefix, not just the completed selector $J$. Insert
inherited eligibility into $J(S)$; it is not a frozen control on an
inner bit shear whose target can change $S$ temporarily. Individual
inner shears still read only the other bit plane and the selector bank.
Their completed transposition preserves $S$. The comparison with the
full intended route is made only after a \emph{complete logical}
mask-selected transposition, source-conditioned operation (A.12), SUM
or recursive basis placement on the first $e$ active ranks, with all
its dirty selectors restored. Write this
prefix implementation as $P_u$ and its intended permutation as $T$.
On a good eligible prefix all intended endpoints are inside it, so
$P_u=T$. On a bad or ineligible prefix the selector $J$ is zero:
the two raw transpositions in (A.11) cancel, each inner operation is
identity, and the two source-predicate toggles in (A.12) cancel too.
Thus the completed $P_u$ is identity there. These statements do not
require raw dirty $T_Z$ to preserve paired sources.

One must not repair an internal truncated $T_Z$ to its full-$h$
version: with $h=2,u=e=1$, $d=(0,0)$ and $Z=(0,1)$, the prefix is
good but full $T_Z$ for $(0\ 1)$ changes the second digit and truncated
$T_Z$ does not. This toy example embeds for any actual $u<h$ by
putting the dirty selector immediately after the prefix and
using active digits throughout the prefix.
That off-prefix action is intended to cancel in the
outer word. Nor is the prefix bad set invariant after an inner bit
shear: the first shear of $(1\ 2)$ sends idle $(X,Y)=(1,1)$ to active
$(0,1)$. Its tile repair uses a different invariant, as specified in
the native appendix.
The separate top-level digit-group placement can reach ranks beyond
$e$; it retains its full-$h$ root-level charge and is not justified
by this width-$e$ prefix argument.
In the uniform $[4]^h$
address space the prefix active count is
binomial with mean $3u/4$. For $u\ge2e$ and $u<h$, a Chernoff bound gives
\[
 \Pr\{B_u<e\}\le e^{-u/24}.                            \tag{A.15}
\]
For $u=h$, eligibility $|S|\ge e$ leaves no exceptional eligible mask.
Both $T$ and the completed logical $P_u$ preserve $S$, so the
exceptional set
$E_{\rm inherited}(S)\wedge[|S|\ge e]\wedge[B_u<e]$ is invariant.
Its whole-universe fraction is bounded by the unconditional binomial
estimate; no conditional distribution given eligibility is assumed.
At that completed boundary, sorting the exceptional records by
$T(P_u^{-1}(y))$ and refilling their slots repairs it exactly;
the maps agree everywhere else. No sector is dropped. The full role recursion always
stores the complete $[4]^h$ address universe, and its row split depends
only on the independent role field, never on $S$. Thus every child
retains the same mask distribution and off-eligibility maps are identity.
Consequently, on every sector $S$, the role program applies
$C^{\otimes |S|}$ to its active digits and identity to its idle digits;
the direct sum over all masks is exactly $A^{\otimes h}$.

The prefix correction above belongs to the completed logical routing
permutation, including $T$ and $T^{-1}$ below, after all its source
conditioning and before a linear combination forms a phase. An
internal $T_{J\wedge Q}$ is not separately corrected when its surrounding
source predicate was prefix-truncated too. It never sorts the output
of the non-permutation matrix $C$.

At selected width one, let $T$ cyclically permute that eligible active
trit and be identity on ineligible masks. The mask-only transpositions
above implement $T$ and $T^{-1}$. Three temporary array copies and a
pointwise linear combination compute
$I+\eta^2(T+T^{-1})$, which is $C$ on the selected trit. Off eligibility
the two identical terms cancel, leaving the identity. The copies,
returns and cleanup add $O(V)$ at this base call. Norm two uses its
inverse via the fixed coefficient wrapper described below.

The $M=49^3$-slot basis changes in the finite frame program decompose
into a fixed number of row-addition matching batches on the first $e$
selected ranks. Nondegenerate residual spaces over $\mathbb F_3$ have
orthogonal bases: repeatedly split a nonisotropic line; its orthogonal
complement remains nondegenerate. The sign of the directed frame
difference must also be retained.

\begin{lemma}[Signed direction and one-child implementation]
\label{lem:signed-phase-edge}
Let $v$ be one such direction in standard-dot coordinates, with
$\nu=\langle v,v\rangle\in\{1,2\}$, and let
$\varepsilon\in\{1,2\}$ specify the positive or negative directed
phase. The effective child coefficient is
\[
             \lambda=\varepsilon/\nu\quad\hbox{in }\F_3.
\]
For the block direction $\langle v\rangle\otimes\F_3^f$, use one
ordinary width-$f$ child if $\lambda=1$, and one conjugated inverse
child if $\lambda=2$. Both cases retain exactly the same child volume.
\end{lemma}
\begin{proof}
Choose $A\in\operatorname{SL}_M(\F_3)$ with $Av=e_1$.
Extend $v$ to a basis and adjust the sign of another column to obtain
determinant one; $M\ge2$. A fixed elementary row-addition word realizes
$A$. Define the address pushforward by
$(Q_Af)(y)=f(A^{-1}y)$. Its translation conjugacy is
$Q_A^{-1}T_wQ_A=T_{A^{-1}w}$.
Chronologically place the addresses by $A\otimes I_f$, apply the
first-block child $(C^{\otimes f})^\lambda$, and undo the placement.
The resulting operator is
\[
 Q_{A\otimes I_f}^{-1}(C^{\otimes f})^\lambda Q_{A\otimes I_f}.
\]
Its Fourier multiplier is
\[
 \eta^{\lambda\sum_{\ell=1}^f\langle v,z^{(\ell)}\rangle^2}.
\]
The desired directed line phase has exponent
\[
 \varepsilon\sum_\ell\langle v,z^{(\ell)}\rangle^2/\nu,
\]
so the stated $\lambda$ is exact. The placement need not preserve the
quadratic form: $\nu$ is computed before placement, in the original
form. Replacing $v$ by $-v$ leaves these squares unchanged.

For $\lambda=2$, wrap the \emph{same forward child} with the
coefficientwise $\F_4$ Frobenius on both existing codec planes,
as in (A.17). This gives $(C^{\otimes f})^{-1}$ on arbitrary codec
values. The wrapper is $\F_2$-linear; its completed conjugate is
linear over the codec field. It adds two linear child-stream scans,
and no second recursive call or new coordinate. If eligibility is
false, the child is identity and the two involutions cancel there.
\end{proof}

Explicitly, the forward/inverse cases are
\[
\begin{array}{c|c|c|c}
 \varepsilon&\nu&\lambda&\hbox{child}\\ \hline
 1&1&1&\hbox{forward}\\
 1&2&2&\hbox{inverse}\\
 2&1&2&\hbox{inverse}\\
 2&2&1&\hbox{forward}
\end{array}
\]
An upward nested-frame edge has $\varepsilon=1$; a downward edge,
including $t_3\to t_4$, has $\varepsilon=2$ on its residual subspace.
The exceptional source difference splits as a positive residual
projection plus $2P_{U_a}$, so its extra unit-line child in (A.9)
also has $\lambda=2$. This flag selects an operator, not an extra
child; the rank sum $s$ is unchanged.

At selected width $e=Mf$, each rank-$a$ edge uses exactly $a$ width-$f$
children, one for each orthogonal basis direction. With $s$ the total
child count (A.9), and $W$ roles, normalized time obeys
\[
 F(e;h)\le (s/W)F(e/M;h)
       +O\!\left(\operatorname{polylog}(L+2)
          (e^\tau+\operatorname{polylog}(L+2))\right).       \tag{A.16}
\]
The additive term includes dirty-selector shears, prefix exception
repairs, basis placements and current-parent tape parking. Put
$\sigma=\log_M(s/W)$. Exact unrolling gives
\[
 O(e^{\max\{\sigma,\tau\}}\operatorname{polylog}(L+2)).
\]
Equality $\sigma=\tau$ adds only $\log e=O(\log L)$. Here the
reconstructed interchange permits $\tau=1-2^{-36}$ and $\sigma<\tau$,
so the ratio $(s/W)/M^\tau$ is strictly below one and the bound is
\[
 O(e^\tau\operatorname{polylog}(L+2)).
\]
No additional exponent
slack is required at this step.
Arbitrary active counts are
partitioned into base-$M$ power-width groups; digit-conditioned groups
partition each mask's active ranks, and each top-level placement is
performed once per group, not again in every child. Hence the complete
$A^{\otimes h}$ tensor has a strict exponent in $h$, relative to the
stated interchange and weighted-mask suppliers.

\subsection{Codec field and explicit positive word for $H_2\otimes H_2$}

The address field $\mathbb F_3$ above is not the phase field. For the
constant word use
\[
 E=\mathbb F_2[t]/(t^4+t+1),\qquad \eta=t^5=t^2+t.
\]
The quartic has no root in $\mathbb F_2$ and remainder $1$ modulo the
only irreducible quadratic $t^2+t+1$, so it is irreducible. Thus
$E=\mathbb F_{16}$, it contains $\mathbb F_4=\mathbb F_2(\eta)$, and
$\eta^2+\eta+1=0$. The four-state phase letter is
\[
 A=(\eta I_3+\eta^2J_3)\oplus[1]=C\oplus[1].
\]
It acts separately on the two $\mathbb F_4$ coefficient planes in
$E=\mathbb F_4\oplus t\mathbb F_4$. Its inverse is $A^2$.
The inverse-phase wrapper on these planes is
\[
 \widehat\sigma(a+tb)=a^2+t b^2\quad(a,b\in\mathbb F_4),       \tag{A.17}
\]
which applies the $\mathbb F_4$ Frobenius to each plane and fixes the
codec basis element $t$. It is $\mathbb F_2$-linear and is not ordinary
$\mathbb F_{16}$ squaring.
In the polynomial basis $(1,t,t^2,t^3)$ it is the explicit map
\[
 (c_0,c_1,c_2,c_3)\longmapsto
 (c_0+c_2+c_3,\ c_1+c_3,\ c_2,\ c_3).
\]
Indeed $a_0=c_0$, $a_1=c_2+c_3$, $b_0=c_1+c_2+c_3$ and
$b_1=c_3$ in $a=a_0+a_1\eta$, $b=b_0+b_1\eta$.
Three in-place XORs implement it, and repeating them gives the identity.
For $\lambda\in\F_4$,
$\widehat\sigma(\lambda z)=\lambda^2\widehat\sigma(z)$.
Consequently coefficientwise conjugation gives
\[
 \widehat\sigma C\widehat\sigma
   =\eta^2I_3+\eta J_3=C^2=C^{-1},
 \qquad
 \widehat\sigma(C^{\otimes f})\widehat\sigma
   =(C^{-1})^{\otimes f}.
\]
This holds on arbitrary $E$-valued inputs. The two passes cancel on
idle or ineligible coordinates; either pass alone can change their
coefficient values. Both passes act within the existing four-bit lanes,
in $O(V)$ time around one child. Ordinary $E$ squaring used on both
sides would fail even on an idle value $t$, sending it to
$t^4=t+1$ instead of $t$. The semilinearity assertion is only for
$\F_4$ scalars, not arbitrary $E$ scalars.

Here is the explicit word construction. Put $a=\eta^2$ and $b=\eta$;
use column vectors and rightmost-first composition. Then
\[
 A^3=I,\quad A^{-1}=A^2,\quad
 u=(1,a,a,0)^T,\quad v=(1,b,b,0)^T,\quad
 P=AE_{11}A^{-1}=uv^T,
\]
with $v^Tu=1$, hence $P^2=P$. Let
\[
 D=\operatorname{diag}(t,1,t+1,1),\quad Q=DPD^{-1}=x_0y_0^T,
\]
where $x_0=(t,a,a(t+1),0)^T$ and
$y_0=(t^{-1},b,b/(t+1),0)^T$. Direct products give
\[
 v^Tx_0=0,\qquad y_0^Tu=\kappa
 =\frac{t^2+t+1}{t(t+1)}\ne0,
 \qquad PQ=0,\quad QP=\kappa x_0v^T.                 \tag{A.18}
\]
For idempotents $P,Q$ with $PQ=0$,
\[
 [I+cP,I+dQ]=I+\frac{cd}{1+c}QP.                    \tag{A.19}
\]
This follows by using $(I+cP)^{-1}=I+c(1+c)^{-1}P$ and
$QPQ=0$. With $c=d=\eta^2$, $cd/(1+c)=1$. Thus
\[
\begin{aligned}
 G&=A\operatorname{diag}(\eta,1,1,1)A^{-1}=I+\eta^2P,\\
 H&=DGD^{-1}=I+\eta^2Q,\\
 T_0&=GHG^{-1}H^{-1}=I+\kappa x_0v^T.
\end{aligned}\tag{A.20}
\]
Each of $G,H,G^{-1},H^{-1}$ uses three forward $A$ letters after
replacing inverse $A$ by $A^2$, so $T_0$ uses 12.

The remaining conjugations isolate an elementary matrix. Set
\[
 B=\operatorname{diag}(1,\eta,\eta,1),\quad
 x=(t,1,t+1,0)^T,\quad w=(1,1,1,0)^T.
\]
Then $T_1=BT_0B^{-1}=I+\kappa xw^T$ (12 $A$ letters). Let $P_{12}$
swap coordinates 1 and 2. Since $w$ is fixed and $w^TP_{12}x=0$,
\[
 T_2=T_1P_{12}T_1P_{12}^{-1}
   =I+\kappa(t+1)(e_1+e_2)w^T                       \tag{A.21}
\]
uses 24 $A$ letters. With $D_3=\operatorname{diag}(1,1,\eta,1)$,
\[
 T_3=T_2D_3T_2D_3^{-1}
   =I+\kappa(t+1)\eta(e_1+e_2)e_3^T                \tag{A.22}
\]
uses 48; the two cross pairings vanish and the row sum is
$(1+\eta^{-1})e_3^T=\eta e_3^T$. With
$D_1=\operatorname{diag}(\eta,1,1,1)$,
\[
 T_4=T_3D_1T_3D_1^{-1}
   =I+\kappa(t+1)E_{13}                             \tag{A.23}
\]
uses 96; the cross pairings vanish and the column factor is
$1+\eta=\eta^2$. Put $\mu=\kappa(t+1)\ne0$; then
\[
 T_5=\operatorname{diag}(\mu^{-1},1,1,1)
       T_4\operatorname{diag}(\mu,1,1,1)=I+E_{13}.  \tag{A.24}
\]
All conjugating factors are invertible monomial matrices. Coordinate
permutations produce every $U_{ij}=I+E_{ij}$ with the same 96 forward
$A$ letters. This is an explicit formula, not an uncomputed word search.

Finally write $H_2=US$, with $U=I+E_{12}$ and $S$ the two-coordinate
swap. In ordered coordinates $(00,01,10,11)$, put
\[
 P_{\rm first}=(1\ 3)(2\ 4),\qquad
 P_{\rm second}=(1\ 2)(3\ 4).
\]
Then, using the four explicit elementary-matrix words,
\[
 H_2\otimes H_2
 =U_{13}U_{24}P_{\rm first}U_{12}U_{34}P_{\rm second}.       \tag{A.25}
\]
There are four $U_{ij}$ factors, hence $4\cdot96=384$ forward $A$
letters; the remaining factors expand to 392 invertible monomial
letters, for 776 total. The completed matrix has entries in $\mathbb F_2$.
As an $\mathbb F_2$-linear map on the four binary coordinates of an
$E$-symbol, it applies $H_2\otimes H_2$ independently to each coordinate
plane. Thus a binary $0/1$ embedding returns to $0/1$, and any zero lane
used for padding returns to zero at the completed word. An arbitrary
$E$-valued input need not return to the scalar subfield. Intermediate
$E$ lanes must not be cropped.
When four arbitrary payload bits are packed into one $E$ symbol,
all four coefficient planes remain data; only extra zero bits from
the final partial group are removed. No projection $E\to\F_2$
is applied to such a packed symbol.

Pairing $2h$ identical native endpoint axes and tensoring this fixed word
reduces their algebraic action to $A^{\otimes h}$ letters and fixed
monomial letters. The permutation part of each monomial is in
$\mathrm{AGL}(2,\mathbb F_2)=S_4$: its NOT/CNOT address shears and the
product of its $h$ diagonal $E$ scalars are charged by the separately
supplied native mask/routing and record-scan interfaces. The active-sector
role construction reduces the $A^{\otimes h}$ letters. The native
large-even-radix endpoint appendix supplies their tape price. No binary
$2^k$-axis H2 bound is asserted here; that requires a distinct binary
mask-routing argument and is not needed for the manuscript's native
endpoint interface.
For precision, write a monomial as $M_{\pi(d),d}=m_d$, with
$m_d\in E^*$. Apply the coefficient $\prod_i m_{d_i}$ at its current
\emph{source} address, then route by $\pi^{\otimes h}$. Equivalently
an output-address scan would use $\prod_i m_{\pi^{-1}(e_i)}$.
Using the source-indexed list unchanged after routing is generally
incorrect. Tensoring the fixed word is letterwise:
$(L_1\cdots L_{776})^{\otimes h}
 =L_1^{\otimes h}\cdots L_{776}^{\otimes h}$.
Thus there are 384 tensor phase calls, not 384 separate calls per axis.

\paragraph{Scope.}
The mask shears used here are constructed independently in
Section~\ref{sec:weighted-native-mask}; interchange is proved in
Appendix~\ref{app:swap}. The mask lemma never calls the later native
H2 tensor theorem. Conversely the rational interchange proof uses
only the scalar incidence motif and finite subspace inclusions, not
this phase recursion. These cross-references follow an acyclic
dependency order, despite the order in which appendices are printed.

Equations (A.5)--(A.10) give the finite scratch-safe routing and strict
rank saving; (A.11)--(A.16) give the active/idle and dirty-selector
interface with inherited mask eligibility; (A.17)--(A.25) give the
explicit positive word. Their tape implementation is relative to the
separately proved weighted-native-mask and equal-chunk interchange
suppliers. The appendix does not reprove those tape suppliers, the ordinary
polynomial multiplication composition, or external acceptance of the
claimed asymptotic result.

